\label{sec:arizona}

\subsection{Background: The AIRC Structure}

The Arizona Independent Redistricting Commission was established by Proposition
106, a voter initiative adopted in 2000, following decades of legislative
gerrymandering.  The AIRC consists of five members: two Democrats, two
Republicans, and one independent chair selected by the other four members from
a pool maintained by the state commission on appellate court appointments.  The
Commission drew the 2002, 2012, and 2022 congressional maps.

The U.S. Supreme Court upheld the AIRC's constitutional authority in
\textit{Arizona State Legislature v.\ Arizona Independent Redistricting
Commission}~\citep{airc2015}, holding that the Elections Clause's reference to
``the Legislature'' encompasses a lawmaking body created by the People through
the initiative process.  The AIRC is therefore a constitutionally valid
redistricting body for federal congressional maps.

Arizona's redistricting criteria, in statutory priority order, are:
(1) population equality; (2) VRA compliance; (3) geographic compactness
and contiguity; (4) preservation of communities of interest; and (5)
partisan fairness.  The algorithm addresses criteria 1 and 3 directly, and
criterion 2 through its VRASection weight; criteria 4 and 5 are the domain
of Commission judgment.

\subsection{How the AIRC Would Integrate Algorithm Output}

Under the Commission pathway, algorithmic redistricting functions as a
\emph{constraint-satisfying starting plan}.  The workflow would proceed in
four stages:

\begin{enumerate}
  \item \textbf{Data preparation and algorithm run.}  Commission staff
    (or a designated technical contractor) run \texttt{redist state --state AZ
    --year 2020 --version v1} to generate a population-balanced, compact map.
    The VRASection weight is applied to protect minority-opportunity districts
    in the Phoenix metropolitan area.  The output is a GEOID-to-district
    assignment file and a map visual.

  \item \textbf{Criteria-compliance verification.}  Commission staff run
    \texttt{redist analyze --state AZ --year 2020 --version v1 --types
    compactness demographic} to generate the criteria-compliance report.
    This report documents population balance, compactness scores, and
    minority-population percentages for each district.

  \item \textbf{Public comment period.}  The Commission publishes the
    algorithm output as a draft plan and receives public comment, including
    community-of-interest testimony.  The Commission retains authority to
    adjust boundaries to address specific community-of-interest concerns
    identified in testimony, subject to the constraint that adjustments must
    not reduce population balance below $\pm 1\%$ or reduce compactness
    scores below a floor.

  \item \textbf{Final adoption.}  The Commission votes on the final plan.
    If the algorithm output is adopted without adjustment, the Commission's
    decision record documents that criteria 1--3 are satisfied by the
    algorithm and that criteria 4--5 are satisfied by the absence of
    challenged community separations and by the near-symmetric partisan
    lean of competitive districts.
\end{enumerate}

\subsection{Community-of-Interest Adjustments}

The most delicate question is how the Commission can exercise judgment about
communities of interest without re-introducing political manipulation.  Two
guardrails are essential.

First, community-of-interest adjustments must be \emph{geographically
justified}: they must unite a geographically defined community (a county,
a Native American tribal area, an established municipality) rather than
unite partisan voters who happen to live in dispersed locations.

Second, adjustments must be documented with a \emph{specific factual finding}
that the original algorithm output separated the identified community.
Arizona's statute already requires this; algorithmic redistricting strengthens
the requirement by providing a clean baseline against which adjustments are
measured.

\subsection{Legal Resilience of the Commission Pathway}

The Commission pathway is legally robust for several reasons.  The AIRC's
constitutional authority is established.  Algorithmic redistricting reduces
the decisional surface available to commissioners, reducing opportunities
for the kind of intentional-partisanship finding that would be required to
sustain a challenge.  And the transparency of the algorithm's criteria---
population equality and compactness, publicly documented---makes it difficult
to allege hidden partisan purpose in the starting plan.
